Lack of scalability is a key issue for virtual-environment technology, and more generally for any massive online experience because it prevents the emergence of a truly massive virtual-world infrastructure (Metaverse). The Solipsis project tackles this issue through the use of peer-to-peer technology, and makes it possible to build and manage a world-scale Metaverse in a truly distributed manner. Following a peer-to-peer scheme, entities collaborate to build up a common set of virtual worlds. We describe in this paper the communication protocol used to share data between peers connected together according to a N-Dimensional Voronoi based peer-to-peer network called RayNet. The policy of data exchange between peers takes advantage of the view-depedent representation used for 3D contents. Moreover, the solution proposed to distribute one each node processes usually executed on the server-side such as detection collision and physics computation. Finally, we present our web composant, a 3D navigator that can be easily executed on terminals with low ressources, and that provides solutions to swap between Web 2.0 and Web.

\textbf{Keywords:}
Peer-to-peer System, Metaverse, Shared Virtual Worlds, Massively Decentralized System, Adaptative 3D Streaming.
